% paper5_methods.tex -- Materials and Methods of Paper 5, input by
% papers/paper5_amendments.tex.
\section{Materials and Methods}

\subsection{Treatments and design}
Three organic amendments, cattle manure (MANURE), compost (COMPOST) and
chopped crop residues (RESIDUE), were compared with an unamended control
(CTRL). Each amendment supplied the same total amount of carbon,
\SI{12.0}{\tonne\per\hectare} of carbon over five years, so that differences in
retention reflect the material rather than the rate. The treatments were
arranged in randomised complete block\index{randomised complete block design}s with four blocks at each of the three
sites \parencite{gomez1984}. The data are held in
\texttt{data/trials/t5\_amendments.tsv}.

\subsection{Carbon stocks}
Soil organic carbon was determined by wet oxidation\index{soil organic carbon!wet oxidation} in topsoil sampled to a
depth $d = \SI{0.20}{\metre}$ at the start and end of the trial. Concentrations
were converted to stocks with the bulk density\index{bulk density} of each site,
\begin{equation}
  C = \mathrm{SOC}\times \mathrm{BD}\times d \times 10
  \quad (\si{\tonne\per\hectare}),
  \label{eq:stock}
\end{equation}
with SOC in \si{\gram\per\kilo\gram}, BD in \si{\gram\per\centi\metre\cubed}
and $d$ in metres.

\subsection{Carbon retention}
The change in stock of each plot is
$\Delta C = C_{\mathrm{final}} - C_{\mathrm{initial}}$. The fraction of the
added carbon retained in the soil is calculated against the control of the same
block,
\begin{equation}
  f_{\mathrm{ret}} = \frac{\Delta C_{\mathrm{amended}} - \Delta C_{\mathrm{control}}}{C_{\mathrm{added}}},
  \label{eq:retention}
\end{equation}
so that the slow change in carbon that occurs without inputs is removed. Mean
retention at each site was compared with the clay content of the site, since
protection by association with fine mineral particles is expected to increase
with clay \parencite{six2002,vonlutzow2006}.
